\documentclass{article}
\begin{document}
This research work is focused on detecting low-grade glioma tumorous cells in MRI images. Glioma is a common brain tumor, that exhibits properties of benign tumors[1]. We used the TCGA-LGG Segmentation dataset[2] for our research. It consists of 3929 brain tumor images and corresponding FLAIR abnormality segmentation masks obtained from 110 patients.

Table-1  lists the models used as encoder for U-Net architecture.
 \begin{table}
    \centering
    \caption{Models used for U-Net encoder and trainable blocks/stages for finetuning}
    \begin{tabular}{l c c}
    \hline
    Family & Model & Trainable Blocks \\
    \hline
    EfficientNet & EfficientNetB0 to B7 & Block 30 to 32\\
    DenseNet & DenseNet169, DenseNet201 & Block 7\\
    ResNet &  ResNet18, ResNet50, ResNet101 & Stage 4\\
    \hline
    \end{tabular}
\end{table}

\section*{References}
\begin{enumerate}
    \item [1] A. Wadhwa, A. Bhardwaj, and V. S. Verma, “A review on brain tumor
    segmentation of mri images,” Magnetic resonance imaging, vol. 61, pp. 247-
    259, 2019.
  
    \item [2] M. Buda, A. Saha, and M. A. Mazurowski, “Association of genomic subtypes
    of lower-grade gliomas with shape features automatically extracted by a deep
    learning algorithm,” Computers in biology and medicine, vol. 109, pp. 218-
    225, 2019.
    
    
  
\end{enumerate}


\end{document}